\section{Related Work}\label{sec:related}

\paragraph{Liquidation-based lending.} Pooled lending protocols and collateralized debt positions mark collateral to an oracle price and liquidate below a threshold; the mechanics and their failure modes are surveyed in \citep{bartoletti2021,werner2022sok,gudgeon2020}. Empirical and theoretical work documents that liquidations are procyclical and that leverage transmits shocks \citep{perez2021,qin2021liquidations,heimbach2024leverage,tian2025,lehar2022,chiu2023}. Stablecoin models show how endogenous collateral produces deleveraging spirals \citep{klagesmundt2020,klagesmundt2021,kozhan2021}, and the Terra collapse is the canonical run on an endogenously backed claim \citep{liu2023terra}; \citet{ma2025} study runs on reserve-backed stablecoins when redemption is in the hands of a few arbitrageurs. crvUSD softens liquidation by converting collateral continuously in a dedicated AMM, but it still uses price oracles and retains a hard liquidation when a loan's health falls below zero \citep{crvusd,crvusddocs}. Our setting differs in one structural respect: the collateral's lower value bound is a redemption right against a reserve the lender holds, not a market price.

\paragraph{Lending without an external price.} Removing the external price does not by itself remove liquidation. Ajna has no external price feeds; lenders choose the prices at which they lend, a borrower who is under-collateralized with respect to the lowest utilized price is eligible for liquidation, and a liquidation is started by a kicker who posts a bond, or by a lender \citep{ajna}. Olympus Cooler~V2 lends a stablecoin against gOHM in perpetual positions. It does not use external market prices; governance sets an origination limit per unit of collateral and a higher liquidation limit. Debt grows by interest, and an account whose debt exceeds its collateral times the liquidation limit can be liquidated \citep{nethermindcooler}. Cooler is close to the present class in how it values collateral, and differs in that accrual can carry a position across a liquidation threshold. \Cref{sec:carry} shows how much accrual a floor-anchored loan can carry without one.

\paragraph{Floor-price issuers.} Nirvana combined a virtual AMM, a rising floor and loans at the floor value: a holder locked the token ANA and minted the protocol's own stablecoin NIRV up to the floor value of the lock, with an origination fee and no interest \citep{nirvanadocs}. In July 2022 an attacker used a flash loan to buy ANA, moved its curve price from about \$8 to \$24, and extracted about \$3.5 million from the treasury \citep{certiknirvana}. Baseline pairs a protocol-owned market with a floor (``BLV''). A 2024 security review of Baseline~V2.1 already discusses reserves that leave the floor position on borrowing and are ``accounted for in the virtual reserves'', and analyzes how this interacts with the system's capacity invariant \citep{guardianbaseline}. Baseline's current documentation describes borrowing the full floor value of the locked tokens at zero interest, with a one-time origination fee, no expiry and no liquidation threshold; it states that the floor can never decrease, and that borrowers continue to receive trading fees on their locked tokens \citep{baselinedocs}. Debt-aware bonding curves \vone{} combine a stepped issuance curve with loans at a fraction of the floor value, and give invariant-style arguments for reserve coverage, floor monotonicity and loan safety, together with discussions of leverage, fees and extraction.

Debt-aware reserve accounting and floor-anchored perpetual credit are thus established practice, and the earlier paper contains a reserve-coverage argument. The contributions of the present paper are specific: one induction over a complete set of operations, for arbitrary schedules; a characterization of credit capacity under contractual release rights; the separation of safety from redistribution; conditional results on borrower behavior; a participation result; and exact conditions for two extraction strategies. The safety theorem is an elementary invariant argument. Its value is the specification: a complete list of operations and guards under which the argument closes, stated precisely enough to be machine-checked and to be compared with an implementation line by line.

Three earlier works are particularly close to the floor-cap model, and \cref{tab:novel} states what each already contains and what this paper adds. The mechanism is not new here: lending the reserve against the issuer's own floor-backed tokens, a reserve invariant, locks that stay if a loan is never repaid, and recursive leverage are all in \citet{dabc2026}, and the accounting of lent reserves appears in audits of deployed code \citep{guardianbaseline}. Our question is narrower: under exactly which conditions an issuer can lend its redemption backing without weakening the claims that remain redeemable, and what that safety does and does not imply for the economics.

\begin{table}[t]
\centering\small
\caption{The three closest prior works and what this paper adds to each.}\label{tab:novel}
\begin{tabularx}{\textwidth}{@{}>{\raggedright\arraybackslash}p{0.15\textwidth}>{\raggedright\arraybackslash}X>{\raggedright\arraybackslash}X>{\raggedright\arraybackslash}X@{}}
\toprule
Prior work & Already there & Formal statement new here & Failure exposed or economic implication new here\\
\midrule
\citet{dabc2026} & Floor-anchored loans from the reserve on a stepped curve; a coverage argument; permanent locks; recursive leverage & \Cref{thm:safety} for arbitrary schedules and a complete operation set, machine-checked; future-repayment coverage and its floor boundary (\cref{cr:kernel,cr:buffer}); integer refinement (\cref{prop:int}) & The floor guard on reshapes, (G1); locks rounded up (\cref{prop:int}); abandonment when the sell fee at the floor reaches $1-\gamma(1-\phi)$ (\cref{prop:route}); the differences listed in \cref{app:errata}\\
\citet{kirste2025} & Mint-and-burn market makers with a reserve; curve shapes; mechanisms that raise the backing & The debt side of the same balance sheet; exit payoff and exit distortion of a reconfiguration & A funded reconfiguration can still move value between holders (\cref{sec:reconf})\\
\citet{xiazhou2007} & A collateralized loan as a perpetual American option & Deferral and break-even horizon when the floor is the lender's own liability and the loan rate is zero (\cref{prop:repay,prop:util}) & The loan is a transfer of time value (\cref{prop:issuer}); in the stationary floor economy with $g<\rho$ and $\gamma+f_s<1$ the mean entry value is negative (\cref{thm:entry})\\
\bottomrule
\end{tabularx}
\end{table} \Cref{tab:class} compares the designs with our model. We do not claim that any deployed system is an instance of the model. Whether it is depends on details that \cref{sec:conformance} lists, and two differences are visible already: a loan at the full floor value is the boundary case $\gbar=1$ of \cref{thm:safety}, which the incentive results cover only in the extended form of \cref{sec:general}; and a loan paid out in a token that the issuer itself creates, as in Nirvana, is not a withdrawal from a reserve and needs a model of that token's liabilities.

\begin{table}[t]
\centering\small
\caption{Protocol-native credit designs and the model of this paper. Sources: \citep{nethermindcooler,ajna,nirvanadocs,baselinedocs,dabc2026}.}\label{tab:class}
\begin{tabularx}{\textwidth}{@{}l>{\raggedright\arraybackslash}X>{\raggedright\arraybackslash}X>{\raggedright\arraybackslash}X@{}}
\toprule
Design & Loan limit & Carrying cost & Forced closure\\
\midrule
Ajna & from lender-chosen prices & interest & when under-collateralized against the lowest utilized price\\
Olympus Cooler V2 & governance-set origination limit & interest & when debt exceeds the liquidation limit\\
Nirvana (2022) & floor value, paid in own stablecoin & none; one-time fee & none by design\\
Baseline (2026) & full floor value & none; one-time fee; locked tokens earn fees & none by design\\
DABC \citep{dabc2026} & fraction $\gamma<1$ of floor value & none; one-time fee & none by design\\
Floor-cap model (this paper) & fraction $\gamma\le\gbar\le1$ of floor value & none, or capped (\cref{prop:carry}) & none (\cref{cor:liqfree})\\
\bottomrule
\end{tabularx}
\end{table}

\paragraph{Bonding curves and primary issuance.} Bonding curves as issuance mechanisms with an invariant linking supply and reserve are formalized in \citep{zargham2020}; \citet{kirste2025} analyze supply-sovereign market makers that mint and burn against a reserve. Constant-function market makers for existing inventories are analyzed in \citep{angeris2020uniswap}, and \citet{xu2023sok} survey them; \citet{bartoletti2022amm} give an operational theory of such market makers and \citet{bichuch2026} an axiomatic one. \Citet{chitra2022borrowing} study loans against market-maker shares under liquidation rules based on the value of the collateral, which depends on the pool price; the collateral here is a claim on a floor that the lender itself maintains. Our model takes the curve as an arbitrary non-decreasing price schedule and adds the debt side of the balance sheet.

\paragraph{Aggregation and release rights.}
The adverse residual-debt envelope in \cref{sec:credit} is the classical fractional-knapsack value function \citep{kellerer2004}. Its upper integrated-quantile representation and its ordering with dispersion are established tools \citep{gushchin2017}; equal debt ratios minimize it by a chord comparison. We use these tools to derive the exact reserve requirement associated with independently exercisable collateral-release rights. For whole-position release the analogous condition is over subsets, and \cref{cr:hardness} applies the classical PARTITION reduction \citep{karp1972}. Neither algorithmic ingredient is presented as a new general mathematical method. The contract-specific conclusions are the necessity of the individual floor cap under proportional release, the reserve needed to exceed it, and the withdrawal, splitting and repricing constraints of alternative caps.

\paragraph{Participation and platform finance.}
\Citet{cong2021tokenomics} jointly model token adoption and valuation through transactional demand and network effects; \citet{cong2022finance} study token issuance and investment in platform productivity. \Citet{limann2025} show how costly observable token purchases can coordinate subsequent platform use. Our conditional floor economy does not replace these demand models. Its participation identity locates a missing benefit in the cash-only benchmark, and the elementary deposit comparison in \cref{prop:deposit-benchmark} states the capital commitment that a service must compensate. Whether redeemability and reserve credit improve adoption or weaken a commitment signal remains an economic question outside the proved results.

\paragraph{Runs and options.} The redemption side of the mechanism resembles a demand-deposit contract in the sense of \citet{diamond1983}, with the difference that each redeemer's curve return is fixed by a schedule that the reserve covers in full; we use the comparison as an analogy (\cref{cor:run}). The borrower's claim is a perpetual American call. \Citet{xiazhou2007} introduced this reading for stock loans, in which a borrower may repay and recover the stock or abandon it (\citet{daixu2011} treat the finite-maturity case), and our repayment result is in the spirit of the no-early-exercise argument of \citet{merton1973}. A loan whose only enforcement is the loss of the collateral is also the building block of collateral equilibrium \citep{geanakoplos2014}; there equilibrium selects which contracts are traded from a set that is given, whereas here contracts are evaluated against an issuer-maintained redemption reserve. The floor-cap contract uses the floor anchor; \cref{sec:credit} characterizes how different release rights change its feasible obligations.

\paragraph{Extractable value.} Front-running and sandwich attacks on automated market makers are documented and analyzed in \citep{daian2020,zhou2021hft,qin2022bev,heimbach2022sandwich}; loss-versus-rebalancing quantifies the adverse selection of passive liquidity \citep{milionis2022lvr}, and \citet{capponi2025} study how competition for extractable value distorts the allocation of block space. A mint-and-burn curve has no passive liquidity providers, but it does have a deterministic price path, and \cref{sec:extraction} analyzes the sandwich on that path directly.
